% Part: model-theory
% Chapter: models-of-arithmetic
% Section: computable-models

\documentclass[../../../include/open-logic-section]{subfiles}

\begin{document}

\olfileid[hi]{mod}{mar}{cmp}
\section{अंकगणित के संगणनीय प्रतिरूप}

\begin{explain}
मानक प्रतिरूप~$\Struct{N}$ में दो अच्छे गुण हैं। उसका प्रांत प्राकृत
संख्याएँ~$\Nat$ हैं, अर्थात् उसके अवयव ठीक उसी प्रकार की वस्तुएँ हैं जिनके
बारे में हम अंकगणित की भाषा का उपयोग करके बात करना चाहते हैं, और मानक
संख्यांक~$\num{n}$ वास्तव में~$n$ को ही चुनता है। दूसरा अच्छा गुण यह है कि
~$\Lang{L_A}$ के अतार्किक प्रतीकों की सभी व्याख्याएँ \emph{संगणनीय} हैं।
$\Assign{\prime}{N}$, $\Assign{+}{N}$ और $\Assign{\times}{N}$ के रूप में
कार्य करने वाले परवर्ती, योग और गुणन फलन संख्याओं के संगणनीय फलन हैं।
(मसलन, ट्यूरिंग मशीनों द्वारा संगणनीय, या आदिम पुनरावर्तन द्वारा परिभाष्य।)
और~$\Struct{N}$ पर लघुतर-संबंध, अर्थात् $\Assign{<}{N}$, निर्णेय है।

$\Th{Q}$ और $\Th{PA}$ जैसे अंकगणितीय सिद्धांतों के अमानक प्रतिरूपों में
अमानक अवयव अवश्य होने चाहिए। इसलिए उनके प्रांतों में सामान्यतः~$\Nat$ के
अतिरिक्त !!{element}s भी होते हैं। किंतु किसी भी गणनीय !!{structure} को
~$\Nat$ सहित किसी भी !!{denumerable} समुच्चय पर बनाया जा सकता है। इसलिए
प्रांत~$\Nat$ वाले अमानक प्रतिरूप भी होते हैं। निस्संदेह, ऐसे
प्रतिरूपों~$\Struct{M}$ में कम-से-कम कुछ संख्याएँ अपनी सामान्य भूमिकाएँ नहीं
निभा सकतीं, क्योंकि कोई~$k$,~$\Value{\num{n}}{M}$ से भिन्न होना चाहिए
प्रत्येक~$n \in \Nat$ के लिए।
\end{explain}

\begin{defn}
!!^a{structure}~$\Struct{M}$,~$\Lang{L_A}$ के लिए \emph{संगणनीय} तभी और केवल
तभी है जब $\Domain{M} = \Nat$ हो, $\Assign{\prime}{M}$,
$\Assign{+}{M}$, $\Assign{\times}{M}$ संगणनीय फलन हों और
$\Assign{<}{M}$ निर्णेय संबंध हो।
\end{defn}

\begin{ex}\ollabel{ex:comp-model-q}
\olref[mdq]{ex:model-K-of-Q} की संरचना $\Struct{K}$ याद कीजिए। उसका प्रांत
$\Domain{K} = \Nat
\cup \{a\}$ था और उसकी व्याख्याएँ थीं:
\begin{align*}
  \Assign{\Obj{0}}{K} & = 0\\
  \Assign{\prime}{K}(x) & =
  \begin{cases}
    x+1 & \text{यदि $x\in \Nat$}\\
    a & \text{यदि $x = a$}
  \end{cases}\\
  \Assign{+}{K}(x, y) & =
  \begin{cases}
    x+y & \text{यदि $x$, $y \in\Nat$}\\
    a & \text{अन्यथा}
  \end{cases}\\
  \Assign{\times}{K}(x, y) & =
  \begin{cases}
    xy & \text{यदि $x$, $y \in\Nat$}\\
    0 & \text{यदि $x=0$ या $y=0$}\\
    a & \text{अन्यथा}\\
  \end{cases}\\
  \Assign{<}{K} & =
  \Setabs{\tuple{x,y}}{x, y \in \Nat \text{ और } x<y} \cup
  \Setabs{\tuple{x,a}}{n \in \Domain{K}}
\end{align*}
किंतु $\Domain{K}$ !!{denumerable} है और इसलिए~$\Nat$ के समसंख्य है।
उदाहरणार्थ, वह $g\colon \Nat \to \Domain{K}$ जिसके लिए $g(0) =
a$ और $g(n) = n+1$ जब $n>0$, एक !!a{bijection} है। हम इसे एक नए
प्रतिरूप~$\Struct{K'}$, जो~$\Th{Q}$ का है, और $\Struct{K}$ के बीच समरूपता
बना सकते हैं।
~$\Struct{K'}$ में हमें~$\Lang{L_A}$ के प्रतीकों को भिन्न फलन और संबंध
निर्दिष्ट करने पड़ते हैं, क्योंकि~$\Nat$ के भिन्न !!{element}s मानक और
अमानक संख्याओं की भूमिकाएँ निभाते हैं।

विशेषतः, अब $0$ सबसे छोटी मानक संख्या की नहीं, बल्कि~$a$ की भूमिका निभाता
है। अब सबसे छोटी मानक संख्या~$1$ है। इसलिए हम
$\Assign{\Obj{0}}{K'} = 1$ निर्दिष्ट करते हैं। परवर्ती फलन भी अब भिन्न है:
कोई मानक संख्या, अर्थात् कोई $n > 0$, दिए जाने पर वह अब भी $n+1$ लौटाता है।
किंतु अब $0$,~$a$ की भूमिका निभाता है, जो स्वयं अपना परवर्ती है। इसलिए
$\Assign{\prime}{K'}(0) = 0$। योग और गुणन के लिए इसी प्रकार हमारे पास
\begin{align*}
\Assign{+}{K'}(x, y) & =
  \begin{cases}
    x+y-1 & \text{यदि $x$, $y >0$}\\
    0 & \text{अन्यथा}
  \end{cases}\\
  \Assign{\times}{K'}(x, y) & =
  \begin{cases}
    1 & \text{यदि $x = 1$ या $y = 1$}\\
    xy - x - y + 2 & \text{यदि $x$, $y > 1$}\\
    0 & \text{अन्यथा}\\
  \end{cases}
\end{align*}
और $\tuple{x, y} \in \Assign{<}{K'}$ तभी और केवल तभी जब $x < y$,
$x > 0$ और $y > 0$, अथवा जब $y = 0$।

ये सभी फलन प्राकृत संख्याओं के संगणनीय फलन हैं और $\Assign{<}{K'}$,
~$\Nat$ पर एक निर्णेय संबंध है---किंतु ये~$\Nat$ पर परवर्ती, योग और गुणन
के समान फलन नहीं हैं, और $\Assign{<}{K'}$,~$<$ के समान संबंध नहीं
है~$\Nat$ पर।
\end{ex}

\begin{prob}
कोई ऐसी !!a{structure}~$\Struct{L'}$ दीजिए जिसमें
$\Domain{L'} = \Nat$ हो और जो
\olref[mod][mar][mdq]{ex:model-L-of-Q} की~$\Struct{L}$ की समरूपी हो।
\end{prob}

\begin{explain}
\Olref{ex:comp-model-q} दिखाता है कि~$\Th{Q}$ के प्रांत~$\Nat$ वाले
संगणनीय अमानक प्रतिरूप हैं। किंतु निम्न परिणाम दिखाता है कि~$\Th{PA}$ के
प्रतिरूपों के लिए यह सत्य नहीं है (और इसलिए~$\Th{TA}$ के प्रतिरूपों के लिए
भी नहीं)।
\end{explain}

\begin{thm}[टेनेनबाउम का प्रमेय]
$\Struct{N}$,~$\Th{PA}$ का एकमात्र संगणनीय प्रतिरूप है।
\end{thm}
  
\end{document}
